% ======================================================================
% Chapter 14: Invitation: The Parton Model of Hadron Structure
% (book pages 473-480; PDF pages 496-503)
% ======================================================================
\chapter{Invitation: The Parton Model of Hadron Structure}

In Part~II of this book, we explored the structure of quantum field theories
in a formal way.  We developed sophisticated calculational algorithms
(Chapter~10), derived a formalism for the extraction of scaling laws and
asymptotic behavior (Chapter~12), and worked out some of the consequences of
spontaneously broken symmetry (Chapter~11).  Much of this formalism turned
out to have unexpected applications in statistical mechanics.  However, we
have not yet investigated its implications for elementary particle physics.
To do so, we must first ask which particular quantum field theories describe
the interactions of elementary particles.

Since the mid-1970s, most high-energy physicists have agreed that the
elementary particles that make up matter are a set of fermions, interacting
primarily through the exchange of vector bosons.  The elementary fermions
include the leptons (the electron, its heavy counterparts $\mu$ and $\tau$, and
a neutral, almost massless neutrino corresponding to each of these species),
and the quarks, whose bound states form the particles with nuclear
interactions, mesons and baryons (collectively called hadrons).  These
fermions interact through three forces: the strong, weak, and electromagnetic
interactions.  Of these, the strong interaction is responsible for nuclear
binding and the interactions of the constituents of nuclei, while the weak
interaction is responsible for radioactive beta decay processes.  The
electromagnetic interaction is the familiar Quantum Electrodynamics, coupled
minimally to all charged quarks and leptons.  It is not clear that these three
forces suffice to explain the most subtle properties of the elementary
fermions---we will discuss this question in Chapter~20---but these three
forces are certainly the most prominent.  All three are now understood to be
mediated by the exchange of vector bosons.  The equations describing the
electromagnetic interaction were discovered by Maxwell, and their quantum
mechanical implications have been treated in detail in Part~I.  The correct
theories of the weak and strong interactions were discovered much later.

By the late 1950s, studies of the helicity dependence of weak interaction
cross sections and decay rates had shown that the weak interaction involves a
coupling of vector currents built of quark and lepton fields.%
\footnote{For an overview of weak interaction phenomenology, see Perkins
(1987), Chapter~7, or any other modern particle physics text.}  It was thus
natural to assume that the weak interaction is due to the exchange of very
heavy vector bosons, and indeed, such bosons, the $W$ and $Z$ particles, were
discovered in experiments at CERN in 1982.  But a complete theory of the weak
interaction must include not only the correct couplings of the bosons to
fermions, but also the equations of motion of the boson fields themselves, the
analogue, for the $W$ and $Z$, of Maxwell's equations.  Finding the correct
form of these equations was not straightforward, because Maxwell's equations
prohibit the generation of a mass for the vector particle.  The proper
reconciliation of the generalized Maxwell equations with the nonzero $W$ and
$Z$ masses turned out to require incorporating into the theory a
spontaneously broken symmetry.  Chapters~20 and~21 treat this subject in
some detail, describing the interplay of vector field theories with
spontaneously broken symmetry.  This interplay leads to new twists and new
phenomena, beyond those discussed in our treatment of spontaneous symmetry
breaking in Chapter~11.  A complete theory of the weak interaction also
requires the simultaneous incorporation of the electromagnetic interaction,
forming a unified structure as first hypothesized by Glashow, Weinberg, and
Salam.

On the other hand, it was for a long time completely obscure that a theory of
exchanged vector bosons could correctly describe the strong interaction.
Part of the mystery was that quarks do not exist as isolated species.  Their
existence, and eventually their quantum numbers, had to be deduced from the
spectrum of observable strongly interacting particles.  But, in addition,
there were complications due to the fact that the strong interactions are
strong.  The Feynman diagram expansion assumes that the coupling constant is
small; when the coupling becomes strong, a large number of diagrams are
important (if the series converges at all) and it becomes impossible to pick
out the contributions of the elementary interaction vertices.  The crucial
clue that the strong interactions have a vector character arose from what at
first seemed to be just another mystery, the observation that the strong
interactions turn themselves off when the momentum transfer is large, in a
sense that we will now describe.

\subsection*{Almost Free Partons}

In Section~5.1 we computed the cross section for the QED process
$e^{+}e^{-}\to\mu^{+}\mu^{-}$.  We then remarked that the corresponding cross
section for $e^{+}e^{-}$ annihilation into hadrons could be computed in the
same way, using a simplistic model in which the quarks are treated as
noninteracting fermions.  This method gives a surprisingly accurate formula
for the cross section, capturing its most important qualitative features.
But we deferred the explanation of this puzzle: How can a model of
noninteracting quarks represent the behavior of a force that, under other
circumstances, is extremely strong?

In fact, there are many circumstances in the study of the strong interaction
at high energy in which this force has an unexpectedly weak effect.
Historically, the first of these appeared in proton-proton collisions.  At
high energy, above 10 GeV or so in the center of mass, collisions of protons
(or any other hadrons) produce large numbers of pions.  One might have
imagined that these pions would fill all of the allowed phase space, but, in
fact, they are mainly produced with momenta almost collinear with the
collision axis.  The probability of producing a pion with a large component
of momentum transverse to the collision axis falls off exponentially in the
value of this transverse momentum, suppressing the production substantially
for transverse momenta greater than a few hundred MeV.

This phenomenon of \emph{limited transverse momentum} led to a picture of a
hadron as a loosely bound assemblage of many components.  In this picture, a
proton struck by another proton would be torn into a cloud of pieces.  These
pieces would have momenta roughly collinear with the original momentum of
the proton and would eventually reform into hadrons moving along the
collision axis.  By hypothesis, these pieces could not absorb a large
momentum transfer.  We can characterize this hypothesis mathematically as
follows: In a high-energy collision, the momenta of the two initial hadrons
are almost lightlike.  The shattered pieces of the hadrons, arrayed along the
collision axis, also have lightlike momenta parallel to the original momentum
vectors.  This final state can be produced by exchanging momenta $q$ among
the pieces in such a way that, though the components of $q$ might be large,
the invariant $q^2$ is always small.  The ejection of a hadron at large
transverse momentum would require large (spacelike) $q^2$, but such a process
was very rare.  Thus it was hypothesized that hadrons were loose clouds of
constituents, like jelly, which could not absorb a large $q^2$.

This picture of hadronic structure was put to a crucial test in the late
1960s, in the SLAC-MIT deep inelastic scattering experiments.%
\footnote{For a description of these experiments and their ramifications, see
J.~I.~Friedman, H.~W.~Kendall, and R.~E.~Taylor, \emph{Rev.\ Mod.\ Phys.}
\textbf{63}, 573 (1991).}  In these experiments, a 20 GeV electron beam was
scattered from a hydrogen target, and the scattering rate was measured for
large deflection angles, corresponding to large invariant momentum transfers
from the electron to a proton in the target.  The large momentum transfer was
delivered through the electromagnetic rather than the strong interaction, so
that the amount of momentum delivered could be computed from the momentum of
the scattered electron.  In models in which hadrons were complex and softly
bound, very low scattering rates were expected.

Instead, the SLAC-MIT experiments saw a substantial rate for hard scattering
of electrons from protons.  The total reaction rate was comparable to what
would have been expected if the proton were an elementary particle scattering
according to the simplest expectations from QED.  However, only in rare cases
did a single proton emerge from the scattering process.  The largest part of
the rate came from the \emph{deep inelastic} region of phase space, in which
the electromagnetic impulse shattered the proton and produced a system with
a large number of hadrons.

How could one reconcile the presence of electromagnetic hard scattering
processes with the virtual absence of hard scattering in strong interaction
processes?  To answer this question, Bjorken and Feynman advanced the
following simple model, called the \emph{parton model}: Assume that the
proton is a loosely bound assemblage of a small number of constituents,
called \emph{partons}.  These include quarks (and antiquarks), which are
fermions carrying electric charge, and possibly other neutral species
responsible for their binding.  By assumption, these constituents are
incapable of exchanging large momenta $q^2$ through the strong interactions.
However, the quarks have the electromagnetic interactions of elementary
fermions, so that an electron scattering from a quark can knock it out of
the proton.  The struck quark then exchanges momentum softly with the
remainder of the proton, so that the pieces of the proton materialize as a
jet of hadrons.  The produced hadrons should be collinear with the direction
of the original struck parton.

The parton model, incomplete though it is, imposes a strong constraint on the
cross section for deep inelastic electron scattering.  To derive this
constraint, consider first the cross section for electron scattering from a
single constituent quark.  We discussed the related process of electron-muon
scattering in Section~5.4, and we can borrow that result.  Since we imagine
the reaction to occur at very high energy, we will ignore all masses.  The
square of the invariant matrix element in the massless limit is written in a
simple form in Eq.~(5.71):
\begin{equation}
  \frac{1}{4}\sum_{\text{spins}}|\mathcal{M}|^2
  = \frac{4\pi^2\alpha^2 Q_i^2}{s^2}\ \frac{s^2+u^2}{t^2} ,
  \label{eq:14.1}
\end{equation}
where $s$, $t$, $u$ are the Mandelstam variables for the electron-quark
collision and $Q_i$ is the electric charge of the quark in units of $|e|$.
Recall from Eq.~\eqref{eq:5.73} that, for a collision involving massless
particles, $s + t + u = 0$.  Then the differential cross section in the
center of mass system is
\begin{equation}
  \frac{d\sigma}{dt} = \frac{2\pi\alpha^2 Q_i^2}{s^2}\ \frac{s^2+u^2}{t^2}
  .
  \label{eq:14.2}
\end{equation}
Or, since $t = -s(1-\cos\theta_{cm})/2$,
\begin{equation}
  \frac{d\sigma}{d\Omega_{cm}} = \frac{\alpha^2 Q_i^2}{4s}\ \frac{1 +
    \cos^2\theta_{cm}}{\sin^4(\theta_{cm}/2)}
  \ \cdots .
  \label{eq:14.3}
\end{equation}

\begin{figure}[t]
  \centering
  \includegraphics[width=0.55\textwidth]{fig14.1.png}
  \caption{Kinematics of deep inelastic electron scattering in the parton
    model.}
  \label{fig:14.1}
\end{figure}

To make use of this result, we must relate the invariants $s$ and $t$ to
experimental observables of electron-proton inelastic scattering.  The
kinematic variables are shown in Fig.~\ref{fig:14.1}.  The momentum transfer
$q$ from the electron can be measured by measuring the final momentum and
energy of the electron, without using any information from the hadronic
products.  Since $q^\mu$ is a spacelike vector, one conventionally expresses
its invariant square in terms of a positive quantity $Q$, with
\begin{equation}
  Q^2 = -q^2 .
  \label{eq:14.4}
\end{equation}
Then the invariant $t$ is simply $-Q^2$.

Expressing $s$ in terms of measurable quantities is more difficult.  If the
collision is viewed from the electron-proton center of mass frame, and we
visualize the proton as a loosely bound collection of partons (and continue
to ignore masses), we can characterize a given parton by the fraction of the
proton's total momentum that it carries.  We denote this longitudinal
fraction by the parameter $\xi$, with $0 < \xi < 1$.  For each species $i$ of
parton, for example, up-type quarks with electric charge $Q_i = +2/3$, there
will be a function $f_i(\xi)$ that expresses the probability that the proton
contains a parton of type $i$ and longitudinal fraction $\xi$.  The
expression for the total cross section for electron-proton inelastic
scattering will contain an integral over the value of $\xi$ for the struck
parton.  The momentum vector of the parton is then $p = \xi P$, where $P$ is
the total momentum of the proton.  Thus, if $k$ is the initial electron
momentum,
\begin{equation}
  s = (p+k)^2 = 2p\cdot k = 2\xi P\cdot k = \xi s ,
  \label{eq:14.5}
\end{equation}
where $s$ is the square of the electron-proton center of mass energy.

Remarkably, $\xi$ can also be determined from measurements of only the
electron momentum, if one makes the assumption that the electron-parton
scattering is elastic.  Since the scattered parton has a mass small compared
to $s$ and $Q^2$,
\begin{equation}
  0 \approx (p+q)^2 = 2p\cdot q + q^2 = 2\xi P\cdot q - Q^2 .
  \label{eq:14.6}
\end{equation}
Thus
\begin{equation}
  \xi = x,
  \qquad\text{where}\qquad
  x = \frac{Q^2}{2P\cdot q} .
  \label{eq:14.7}
\end{equation}
From each scattered electron, one can determine the values of $Q^2$ and $x$
for the scattering process.  The parton model then predicts the event
distribution in the $x$--$Q^2$ plane.  Using the parton distribution functions
$f_i(\xi)$, evaluated at $\xi = x$, and the cross-section formula
\eqref{eq:14.3}, we find the distribution
\begin{equation}
  \frac{d^2\sigma}{dx\,dQ^2}
  = \frac{2\pi\alpha^2}{Q^4}\ \sum_i Q_i^2\,f_i(x)\ \Bigl[ 1 + \Bigl(
    1 - \frac{Q^2}{xs} \Bigr)^2 \Bigr] .
  \label{eq:14.8}
\end{equation}
The distribution functions $f_i(x)$ depend on the details of the structure of
the proton and it is not known how to compute them from first principles.
But formula \eqref{eq:14.8} still makes a striking prediction, that the deep
inelastic scattering cross section, when divided by the factor
\begin{equation}
  \frac{2\pi\alpha^2}{Q^4}\ \Bigl[ 1 + \Bigl( 1 - \frac{Q^2}{xs}
    \Bigr)^2 \Bigr]
  \label{eq:14.9}
\end{equation}
to remove the kinematic dependence of the QED cross section, gives a quantity
that depends only on $x$ and is independent of $Q^2$.  This behavior is known
as \emph{Bjorken scaling}.  Indeed, the data from the SLAC-MIT experiment
exhibited Bjorken scaling to about 10\% accuracy for values of $Q$ above
1 GeV, as shown in Fig.~\ref{fig:14.2}.

\begin{figure}[t]
  \centering
  \includegraphics[width=0.6\textwidth]{fig14.2.png}
  \caption{Test of Bjorken scaling using the $e^{-}p$ deep inelastic
    scattering cross sections measured by the SLAC-MIT experiment,
    J.~S.~Poucher, et.\ al., \emph{Phys.\ Rev.\ Lett.} \textbf{32}, 118
    (1974).  We plot $d^2\sigma/dx\,dQ^2$ divided by the factor \eqref{eq:14.9}
    against $x$, for the various initial electron energies and scattering
    angles indicated.  The data span the range
    $1\ \mathrm{GeV}^2 < Q^2 < 8\ \mathrm{GeV}^2$.}
  \label{fig:14.2}
\end{figure}

Bjorken scaling is, essentially, the statement that the structure of the
proton looks the same to an electromagnetic probe no matter how hard the
proton is struck.  In the frame of the proton, the energy of the exchanged
virtual photon is
\begin{equation}
  q^0 \approx \frac{1}{2m}\bigl( Q^2 + \cdots \bigr) ,
\end{equation}
where $m$ is the proton mass.  The reciprocal of this energy transfer is,
roughly, the duration of the scattering process as seen by the components of
the proton.  This time should be compared to the reciprocal of the proton
mass, which is the characteristic time over which the partons interact.  The
deep inelastic regime occurs when $q^0 \gg m$, that is, when the scattering is
very rapid compared to the normal time scales of the proton.  Bjorken scaling
implies that, during such a rapid scattering process, interactions among the
constituents of the proton can be ignored.  We might imagine that the partons
are approximately free particles over the very short times scales
corresponding to energy transfers of a GeV or more, though they have strong
interactions on longer time scales.

\subsection*{Asymptotically Free Partons}

The picture of the proton structure implied by Bjorken scaling was
beautifully simple, but it raised new, fundamental questions.  In quantum
field theory, fermions interact by exchanging virtual particles.  These
virtual particles can
\begin{equation}
  \text{(increase the effective strength of the coupling at short
    distances)} .
\end{equation}
But the analysis of the deep inelastic scattering data suggested that the
strong interaction becomes \emph{weaker} at short distances.  In the early
1970s, this apparent contradiction was resolved by the discovery of
\emph{asymptotic freedom} by Gross, Politzer, and Wilczek.  They showed that
non-Abelian gauge theories---theories of interacting vector bosons---have the
remarkable property that the effective coupling constant tends to zero at
large momenta.  This is precisely the behavior required to make sense of
Bjorken scaling in a quantum field theory of the strong interactions in four
dimensions.  This discovery gave the crucial clue for the construction of the
fundamental theory of the strong interactions.  Apparently, the quarks are
bound together by interacting vector bosons (called \emph{gluons}) of
precisely this type.

However, these gauge theories cannot precisely reproduce the expectations of
strict Bjorken scaling.  The differences between the free parton model and
the quantum field theory model with asymptotic freedom appear when one moves
to a higher level of accuracy in measurements of deep inelastic scattering
and other strong interaction processes involving large momentum transfer.  In
an asymptotically free quantum field theory, the coupling constant is still
nonzero at any finite momentum transfer.  In fact, the final evolution of the
coupling to zero is very slow, logarithmic in momentum.  Thus, at some level,
one must find small corrections to Bjorken scaling, associated with the
exchange or emission of high-momentum gluons.  Similarly, the other
qualitative simplifications of hadron physics at high momentum transfer---for
example, the phenomenon of limited transverse momentum in hadron-hadron
collisions---should be only approximate, receiving corrections due to gluon
exchange and emission.  Thus the predictions of an asymptotically free theory
of the strong interaction are twofold.  On one hand, such a theory predicts
qualitative simplifications of behavior at high momentum.  But, on the other
hand, such a theory predicts a specific pattern of corrections to this
behavior.

In fact, particle physics experiments of the 1970s revealed precisely this
picture.  Bjorken scaling was found to be only an approximate relation,
showing violations that correspond to a slow evolution of the parton
distributions $f_i(x)$ over a logarithmic scale in $Q^2$.  The rate of
particle production in hadron-hadron collisions was found to decrease only as
a power rather than exponentially at very large values of the transverse
momentum, and the particles produced at large transverse momentum were shown
to be associated with jets of hadrons created by the soft evolution of a
hard-scattered quark or gluon.  Most remarkably, the forms of the cross
sections found for these and other deviations from scaling did, finally, give
direct evidence for the vector character of the elementary field that
mediates the strong interaction.

We will review all of these phenomena in Chapter~17, as we study the
particular gauge theory that describes the strong interactions.  First,
however, we must learn how to construct non-Abelian gauge theories and how to
work out their predictions using Feynman diagrams.  Throughout our analysis
of these theories, the renormalization group will play an essential role.  One
of the very beautiful aspects of the study of non-Abelian gauge theories is
the way in which the most powerful general ideas of quantum field theory
acquire even more strength as they intertwine with the specific features of
these particular, intricately built models.  This interplay between general
principles and the specific features of gauge theories will be the major
theme of Part~III of this book.
